\documentclass[11pt]{article}
\usepackage{explorer}
\newcommand{\ExplorerLanguage}{FR}
\newcommand{\ExplorerNumber}{10}
\newcommand{\ExplorerKind}{Problème de synthèse}
\begin{document}
\ExplorerTitle{Première--Terminale · suites}{Sommes télescopiques}
\ExplorerSection{Chercher, justifier, conclure}
Soit \(T_{n}=\sum  \frac{1}{k(k+2)}\) pour k de 1 à n. Établir une expression exacte en fonction de n, puis calculer sa limite. Trouver le plus petit n tel que l’écart à la limite soit strictement inférieur à 0,01.\par
\vspace{5pt}
\end{document}
